
The statistical origin of Bose--Einstein condensation is Bose's counting of
light quanta and Einstein's extension to material particles
\citep{Bose1924,Einstein1924IdealGas,Einstein1925IdealGas}.
For interacting particles, Penrose and Onsager formulated condensation through
a macroscopic eigenvalue of the one-particle density matrix
\citep{PenroseOnsager1956}.  In a translation-invariant gas, the constant
orbital is the natural candidate.  Bogoliubov's theory expands about its
macroscopic occupation to obtain the low-energy excitation spectrum
\citep{Bogoliubov1947}; the hard-sphere calculation of Lee, Huang and Yang
predicts the leading energy correction and depletion
\citep{LeeHuangYang1957}.  Establishing the occupation for the exact equilibrium
state is a separate mathematical question from deriving these expansions.

For hard spheres, rigorous energy bounds began with Dyson
\citep{Dyson1957}, and Lieb and Yngvason obtained the leading dilute
asymptotics \citep{LiebYngvason1998}.  Fournais and Solovej proved the
Lee--Huang--Yang lower bound for interactions including hard cores
\citep{FournaisSolovej2023}; Basti, Brooks, Cenatiempo, Olgiati and Schlein
proved the matching upper bound for hard spheres and then for general
nonnegative short-range interactions
\citep{BastiEtAl2026HardSpheres,BastiEtAl2026GeneralPotentials}.
At positive temperature, Fournais, Junge, Girardot, Morin, Olivieri and Triay
obtained a low-temperature free-energy lower bound that includes hard cores
\citep{FournaisEtAl2026FreeEnergy}.  These results control thermodynamic
energies.  They do not by themselves give a volume-uniform lower bound on a
single orbital's occupation: the one-particle kinetic gap tends to zero as
the box grows.

Rigorous continuum condensation has been established in several scaling
regimes.  Lieb and Seiringer proved complete ground-state condensation in the
trapped Gross--Pitaevskii limit \citep{LiebSeiringer2002}.  At positive
temperature, Deuchert, Seiringer and Yngvason treated a trapped gas in a
combined thermodynamic and Gross--Pitaevskii limit
\citep{DeuchertSeiringerYngvason2019}, and Deuchert and Seiringer treated the
homogeneous canonical gas in Gross--Pitaevskii scaling
\citep{DeuchertSeiringer2020GP}.  The latter theorem allows hard cores, but
their scattering length is proportional to $L/N$ and therefore vanishes at
fixed density as $N$ grows.  Further results of Fournais, Chong, Liang and Nam,
and Junge allow increasing box sizes coupled to increasing dilution
\citep{Fournais2021,ChongLiangNam2026,Junge2026}.  Such coupled limits do not
permit unrestricted volume growth at a fixed positive gas parameter.
Deuchert and Seiringer also established positive-temperature condensation in
a combined semiclassical and mean-field limit for a trapped gas
\citep{DeuchertSeiringer2021MeanField}.

Fixed-parameter thermodynamic condensation also has precedents with different
interactions and state formulations.  Aizenman, Lieb, Seiringer, Solovej and
Yngvason proved low-temperature condensation for a hard-core lattice gas at
half filling \citep{AizenmanLiebSeiringerSolovejYngvason2004}.  S\"ut\H{o}'s
continuum theorem assumes both the potential and its Fourier transform to be
nonnegative, with integrability and regularity hypotheses that exclude a
literal hard core \citep[Theorem~1.2]{Suto2023TwoStep}.  Galanda and Pinamonti
construct algebraic equilibrium states for smooth interactions in coupled
weak-interaction limits and, after renormalization, in Gross--Pitaevskii
scaling with vanishing temperature
\citep{GalandaPinamonti2025Equilibrium,GalandaPinamonti2026GP}.
The present theorem concerns the exact canonical hard-sphere state with the
exclusion distance, positive density, and positive temperature all fixed
before the volume grows.  It gives a positive condensate fraction at
sufficiently low temperature; it does not determine the transition
temperature or the leading depletion.

Two companion articles provide the immediate starting point.
\emph{Ground-state condensation in the dilute hard-sphere gas}
\citep{HS} treats the fixed-parameter ground-state problem and develops
hard-core bridge replacements, movable groups of particles, routes avoiding
unusable cells, and cancellation of conditional normalizers.
\emph{A density-uniform condensate bound for dilute Bose gases}
\citep{B} treats bounded finite-range repulsive potentials at positive
temperature and supplies localization and short-slab trajectory estimates,
as well as the underlying lattice-path construction.  We reproduce the
arguments needed here.  The additional issue is to obtain joint spatial and
trajectory estimates for the hard-core canonical Gibbs state, and to make
the short-path comparison control its one-particle density matrix.

The local analysis builds on established tools.  The hard-core ray estimate
belongs to the Dyson local-energy method, and the subdivision and
two-particle bounds are related to the work of Lieb and Yngvason and of
Lundholm, Portmann and Solovej
\citep{Dyson1957,LiebYngvason1998,LundholmPortmannSolovej2015}.
Particle insertion has antecedents in the nearest-neighbor trial states of
Dyson and Lieb, Seiringer and Yngvason
\citep{Dyson1957,LiebSeiringerYngvason2000}.
Localization into two particle species uses the Fock-space localization
formalism of Derezi\'nski and G\'erard, its many-body development by Lewin,
and its thermal application by Deuchert, Seiringer and Yngvason
\citep{DerezinskiGerard1999,Lewin2011,DeuchertSeiringerYngvason2019}.
Strong subadditivity and monotonicity of relative entropy provide the
entropy comparisons \citep{LiebRuskai1973Entropy,Lindblad1975}.
The use of positive imaginary-time paths goes back to Feynman's formulation
\citep{Feynman1953AtomicLambda}; Ueltschi's account distinguishes cycle
information from one-particle off-diagonal order
\citep{Ueltschi2006FeynmanCycles}.  Here the overlap identity addresses the
one-particle occupation directly.

The long-distance comparison has a separate probabilistic ancestry.
Benjamini, Pemantle and Peres constructed paths with exponential intersection
tails \citep{BPP1998}; H\"aggstr\"om and Mossel developed a construction using
random biases on blocks at several scales \citep{HaggstromMossel1998}.
The finite-endpoint averaging method of Abbe, Massouli\'e, Montanari, Sly and
Srivastava, and its adaptation by Garban and Spencer, turn intersection
control into long-range information or order
\citep{AbbeMassoulieMontanariSlySrivastava2018,GarbanSpencer2023}.
Averaging path likelihoods before taking their second moments also appears
in trail detection and in the point-reassignment argument of Peres and Sly
\citep{AriasCastroCandesHelgasonZeitouni2008,BergerPeres2013,PeresSly2014}.
Our lattice construction follows the explicit skeleton in \citet[Section~8.1]{B}
and the obstacle detours in \citet[Section~7]{HS}, using Tim\'ar's
boundary-connectivity theorem \citep[Theorem~4]{Timar2013}.
The hard-core comparison retains the exact interaction normalizers; after
averaging two route changes, their cost is controlled by mutual encounters
rather than by the length of either route.
